%%%%%%%%%%%%%%%%%%%%%%%%%%%%%%%%%%%%%%%%%%%%%%%%%
%%%%%%%%%%%% chap: title %%%%%%%%%%%%%%%%%
%%%%%%%%%%%%%%%%%%%%%%%%%%%%%%%%%%%%%%%%%%%%%%%%%

\chapter{State of the art}\label{chapter:chap2}

The initial stage involves investigating existing work in the area of research to establish the foundation for the thesis.

\section{Resources}

Scientific articles can be found in different online databases:
\begin{itemize}
    \item ACM Digital Library
    \item Science Direct
     \item IEEE Explore
    \item Annual Reviews
    \item Springer Link
    \item Wiley InterScience
    \item SCOPUS
    \item  Google Scholar
     \item JSTOR (social sciences, arts \& humanities)
\end{itemize}

Identify communities on the topic (specific workshops on national and international conferences).

Some of the articles from the databases are accessible through university IP.


\section{First Step: Take notes}

When reviewing the literature, record systematic notes, ensuring that the information can be utilized at a later stage:
\begin{itemize}
    \item resource author;
    \item resource name (article title, book title, ...);
    \item subject: what is the paper aim;
    \item addressed problem ; 
    \item methods/methodology used to solve the problem;
    \item used algorithms;
    \item test data (e.g. benchmarks, real cases).
\end{itemize}

For a git resource or an existing application:
\begin{itemize}
    \item URL;
    \item application features;
    \item used technologies.
\end{itemize}